\documentclass{article}
\usepackage{fontspec}
\usepackage[vietnamese]{babel}
\usepackage[fontsize=12pt]{scrextend}
\usepackage[left=2cm, right=2cm, top=2cm, bottom=2cm]{geometry}
\usepackage{setspace}
\setstretch{1.5}
\usepackage{amsmath}
\usepackage{amsfonts}
\usepackage{amssymb}
\usepackage{graphicx}
\usepackage{changepage}
\usepackage{color}
\renewcommand{\thesection}{\arabic{section}.}
\renewcommand{\thesubsection}{\thesection\arabic{subsection}}
\renewcommand{\thesubsubsection}{\thesubsection.\arabic{subsubsection}}

\usepackage{tocloft}

\renewcommand{\cftdotsep}{\cftnodots}
\setlength{\parindent}{0pt}


\begin{document}

\begin{center}
    \textbf{\Large LỜI GIỚI THIỆU}
\end{center}

\vspace{0.5cm}

\noindent Toán học là nền tảng của nhiều ngành khoa học và có vai trò quan trọng trong việc giải quyết các vấn đề của đời sống cũng như kỹ thuật. Trong chương trình Toán phổ thông, lũy thừa và hàm số lũy thừa là những nội dung cơ bản nhưng có nhiều ứng dụng thiết thực trong các lĩnh vực như kinh tế, khoa học tự nhiên, công nghệ và mô hình hóa các hiện tượng thực tế.

Với mong muốn tìm hiểu sâu hơn về bản chất của lũy thừa, các tính chất, hàm số lũy thừa cũng như những ứng dụng của chúng trong thực tiễn, tôi lựa chọn đề tài "Lũy thừa, hàm số lũy thừa và ứng dụng" để nghiên cứu. Thông qua đề tài này, người đọc có thể hệ thống lại các kiến thức trọng tâm, hiểu rõ hơn các tính chất của lũy thừa, phương pháp khảo sát hàm số lũy thừa và vận dụng chúng vào giải quyết một số bài toán cũng như các tình huống thực tế.

Nội dung của tiểu luận được trình bày theo trình tự từ cơ sở lý thuyết đến ứng dụng, bao gồm: khái niệm và tính chất của lũy thừa, hàm số lũy thừa, một số bài toán liên quan, các ứng dụng trong thực tế và phần khảo sát, đánh giá kết quả. Qua đó, đề tài góp phần củng cố kiến thức, phát triển tư duy toán học và nâng cao khả năng vận dụng kiến thức vào học tập cũng như cuộc sống.

Mặc dù đã cố gắng hoàn thiện trong quá trình thực hiện nhưng khó tránh khỏi những hạn chế và thiếu sót. Tôi rất mong nhận được những ý kiến đóng góp của quý thầy cô và bạn đọc để đề tài được hoàn thiện hơn.

\newpage
\begin{center}
    \textbf{\Large MỤC LỤC}
\end{center}

\vspace{0.5cm}


\noindent \textbf{LỜI GIỚI THIỆU} \dotfill 1 \\
\noindent \textbf{1. LŨY THỪA} \dotfill 5 \\
\hspace*{0.5cm} 1.1 Khái niệm lũy thừa  \dotfill 5 \\
\hspace*{1cm} 1.1.1 Lũy thừa với số mũ nguyên  \dotfill 5 \\
\hspace*{1cm} 1.1.2 Phương trình $x^{\alpha} = b$ \dotfill 5 \\
\hspace*{1cm} 1.1.3  Căn bậc $n$ \dotfill 6 \\
\hspace*{1cm} 1.1.4 Lũy thừa với số mũ hữu tỉ \dotfill 10 \\
\hspace*{1cm} 1.1.5 Lũy thừa với số mũ vô tỉ \dotfill 10 \\
\hspace*{0.5cm} 1.2 TÍNH CHẤT CỦA LŨY THỪA VỚI SỐ MŨ THỰC  \dotfill 5 \\
\noindent \textbf{2. HÀM SỐ LŨY THỪA} \dotfill 11 \\
\hspace*{0.5cm} 2.1 KHÁI NIỆM \dotfill 5 \\
\hspace*{0.5cm} 2.2   ĐẠO HÀM CỦA HÀM SỐ LŨY THỪA\dotfill 5 \\
\hspace*{0.5cm} 2.3 KHẢO SÁT HÀM SỐ LŨY THỪA $y=x^{\alpha}$ \dotfill 5 \\
\noindent \textbf{3. MỘT SỐ BÀI TOÁN LIÊN QUAN ĐẾN LŨY THỪA} \dotfill 30 \\
\noindent \textbf{4. MỘT SỐ ỨNG DỤNG THỰC TẾ CỦA HÀM SỐ LŨY THỪA} \dotfill 30 \\
\noindent \textbf{TÀI LIỆU THAM KHẢO} \dotfill 46
\newpage


% nội dung 
\section{LŨY THỪA}
\subsection{KHÁI NIỆM LŨY THỪA}
\subsubsection{Lũy thừa với số mũ nguyên}

\textbf{Định nghĩa 1.1} {Cho n là một số nguyên dương.
Với $a$ là số thực tùy ý, lũy thừa bậc $n$ của $a$ là tích của $n$ thừa số $a$}\\
\centerline{$a^n=\underset{n \text{ lần}}{\underbrace{a.a...a}}$}
Với $a\ne 0$

\centerline{$a^0  = 1$}
\centerline{$a^{-n} = \dfrac{1}{a^n}$}\

Trong biểu thức $a^m$, ta gọi $a$ là cơ số, số nguyên $m$ là số mũ.

\textbf{Chú ý 1.2} {$0^0$ và $0^{−n}$ không có nghĩa.}

Lũy thừa với số nguyên có các tính chất tương tự của lũy thừa với số mũ nguyên dương
 
 \textbf{Ví dụ 1.3} {Tính giá trị của biểu thức}

\vspace{0.3cm}

\centerline{$
A=\left(\frac{1}{3}\right)^{-10}\cdot 27^{-3}
+\left(0.2\right)^{-4}\cdot 25^{-2}
+128^{-1}\cdot\left(\frac{1}{2}\right)^{-9}.
$}

\vspace{0.3cm}

\textcolor{blue}{\textit{Giải}}

\centerline{$A=3^{10}\cdot\frac{1}{27^3}
+\frac{1}{0,2^4}\cdot\frac{1}{25^2}
+\frac{1}{128}\cdot2^9
=3+1+4=8.$}

\vspace{0.3cm}

\textbf{Ví dụ 1.4} {Rút gọn biểu thức}

\centerline{$
B=\left[\frac{a\sqrt{2}}{(1+a^2)^{-1}}-\frac{2\sqrt{2}}{{a}^{-1}}\right]\cdot\frac{a^{-3}}{1-a^{-2}}.
\qquad (a\ne0,\ a\ne\pm1)$}


\textcolor{blue}{\text{Giải. }}&\text{Với } $a\ne0,\ a\ne\pm1$ \text{ ta có}\\
\[
\begin{aligned}
B &= \left[a\sqrt{2}(1+a^2)-2\sqrt{2}a\right]\cdot\frac{1}{a^3(1-a^{-2})}\\
  &= (a\sqrt{2}+a^3\sqrt{2}-2a\sqrt{2})\cdot\frac{1}{a^3-a}\\
  &= a\sqrt{2}(a^2-1)\cdot\frac{1}{a(a^2-1)}\\
  &= \sqrt{2}.
\end{aligned}
\]

\subsubsection{Phương trình $x^{\alpha} = b$}
Dựa vào đồ thị của các hàm số $y = x^3$ và $y = x^4$, hãy biện luận theo $b$ số nghiệm của các
phương trình $x^3 = b$ và $x^4 = b$.

\begin{center}
    \includegraphics[width=0.8\textwidth]{anh1.png}
\end{center}

Đồ thị của hàm số $y = x^{2k+1}$ có dạng tương tự đồ thị hàm số $y = x^3$ và đồ thị hàm số $y = x^{2k}$
có dạng tương tự đồ thị hàm số $y = x^4$. Từ đó ta có kết quả biện luận số nghiệm của phương
trình $x^n = b$ như sau:\\
a) Trường hợp $n$ lẻ :\\
Với mọi số thực $b$, phương trình có nghiệm duy nhất.\\
b) Trường hợp $n$ chẵn:\\
Với $b < 0$, phương trình vô nghiệm;\\
Với $b = 0$, phương trình có một nghiệm $x = 0$;\\
Với $b > 0$, phương trình có hai nghiệm đối nhau
\subsubsection{Căn bậc $n$}
Cho số nguyên dương $n$, phương trình\\
\centerline{$a^b=b$}\\
đưa đến hai bài toán ngược nhau :\\
• Biết tính a, tính b.\\
• Biết tính b, tính a.\\
Bài toán thứ nhất là tính lũy thừa của một số. Bài toán thứ hai dẫn đến khái niệm lấy căn của
một số.\\
\textbf{Định nghĩa 1.5} \textit{Cho số thực $b$ và số nguyên dương $n$ ($n \ge 2$). Số $a$ được gọi là \textbf{\color{blue}căn bậc $n$} của số $b$ nếu $a^n = b$.}

Chẳng hạn, $2$ và $-2$ là các căn bậc $4$ của $16$; $-\dfrac{1}{3}$ là căn bậc $5$ của $-\dfrac{1}{243}$. Từ định nghĩa và kết quả biện luận về số nghiệm của phương trình $x^n = b$, ta có:

Với $n$ lẻ và $b \in \mathbb{R}$: Có duy nhất một căn bậc $n$ của $b$, ký hiệu là $\sqrt[n]{b}$.

Với $n$ chẵn và $\begin{cases} 
b < 0 & \text{Không tồn tại căn bậc } n \text{ của } b; \\[5pt]
b = 0 & \text{Có một căn bậc } n \text{ của } b \text{ là số } 0; \\[5pt]
b > 0 & \text{Có hai căn trái dấu } n \text{ ký hiệu giá trị dương là } \sqrt[n]{b}, \text{còn giá trị âm là } -\sqrt[n]{b}.
\end{cases}$

\textbf{\color{blue}Tính chất của căn bậc $n$}

Từ định nghĩa ta có các tính chất sau:

\[
\sqrt[n]{a} . \sqrt[n]{b} = \sqrt[n]{ab};
\]

\[
\frac{\sqrt[n]{a}}{\sqrt[n]{b}} = \sqrt[n]{\frac{a}{b}};
\]

\[
(\sqrt[n]{a})^m = \sqrt[n]{a^m};
\]

\[
\sqrt[n]{a^n} = \begin{cases} 
a, & \text{khi } n \text{ lẻ} \\[5pt]
|a|, & \text{khi } n \text{ chẵn;} 
\end{cases}
\]

\[
\sqrt[n]{\sqrt[k]{a}} = \sqrt[nk]{a}.
\]

\textbf{Ví dụ 1.6} \textit{Rút gọn các biểu thức:}

\begin{tabular}{@{}p{0.5\textwidth} p{0.5\textwidth}@{}}
a) $\sqrt[5]{4} . \sqrt[5]{-8};$ & b) $\sqrt[3]{3\sqrt{3}}.$
\end{tabular}

\textbf{\color{blue}\textit{Giải.}}

\begin{enumerate}
    \item $\sqrt[5]{4} . \sqrt[5]{-8} = \sqrt[5]{-32} = \sqrt[5]{(-2)^5} = -2.$
    \item $\sqrt[3]{3\sqrt{3}} = \sqrt[3]{(\sqrt{3})^3} = \sqrt{3}.$
\end{enumerate}

\subsubsection{Lũy thừa với số mũ hữu tỉ}
\textbf{Định nghĩa 1.7} \textit{Cho số thực $a$ dương và số hữu tỉ $r = \dfrac{m}{n}$, trong đó $m \in \mathbb{Z}, n \in \mathbb{N}, n \ge 2$. Lũy thừa của $a$ với số mũ $r$ là số $a^r$ xác định bởi}
\[
a^r = a^{\frac{m}{n}} = \sqrt[n]{a^m}.
\]

\textbf{Ví dụ 1.8}
\[
\begin{aligned}
\left(\frac{1}{8}\right)^{\frac{1}{3}} &= \sqrt[3]{\frac{1}{8}} = \frac{1}{2}; \\[8pt]
4^{-\frac{3}{2}} &= \sqrt{4^{-3}} = \frac{1}{\sqrt{4^3}} = \frac{1}{8}; \\[8pt]
a^{\frac{1}{n}} &= \sqrt[n]{a} \quad (a > 0, n \ge 2).
\end{aligned}
\]

\textbf{Ví dụ 1.9} \textbf{Rút gọn biểu thức}
\[
D = \frac{x^{\frac{5}{4}}y + xy^{\frac{5}{4}}}{\sqrt[4]{x} + \sqrt[4]{y}} \qquad (x, y > 0).
\]

\textbf{\color{blue}\textit{Giải.}} Với $x$ và $y$ là những số dương, theo định nghĩa, ta có
\[
D = \frac{xy(x^{\frac{1}{4}} + y^{\frac{1}{4}})}{x^{\frac{1}{4}} + y^{\frac{1}{4}}} = xy.
\]

\subsubsection{Lũy thừa với số mũ vô tỉ}

Cho $a$ là một số dương, $\alpha$ là một số vô tỉ. Ta thừa nhận rằng luôn có một dãy số hữu tỉ $(r_n)$ có giới hạn là $\alpha$ và dãy số tương ứng $(a^{r_n})$ có giới hạn không phụ thuộc vào việc chọn dãy $(r_n)$.

\textbf{Định nghĩa 1.10} \textit{Ta gọi giới hạn của dãy số $(a^{r_n})$ là lũy thừa của $a$ với số mũ $\alpha$, ký hiệu là $a^\alpha$:}
\[
a^\alpha = \lim_{n \to +\infty} a^{r_n} \quad \text{với } \alpha = \lim_{n \to +\infty} r_n.
\]

\textbf{Chú ý 1.11} \textit{Từ định nghĩa, ta có } $1^\alpha = 1 \qquad (\alpha \in \mathbb{R})$.

\
\subsection{TÍNH CHẤT CỦA LŨY THỪA VỚI SỐ MŨ THỰC}

Cho $a, b$ là những số thực dương; $\alpha, \beta$ là những số thực tùy ý. Khi đó, ta có:

\[
\begin{aligned}
a^\alpha . a^\beta &= a^{\alpha+\beta}; \\[4pt]
\frac{a^\alpha}{a^\beta} &= a^{\alpha-\beta}; \\[4pt]
(a^\alpha)^\beta &= a^{\alpha . \beta}; \\[4pt]
(ab)^\alpha &= a^\alpha . b^\alpha; \\[4pt]
\left(\frac{a}{b}\right)^\alpha &= \frac{a^\alpha}{a^\beta};
\end{aligned}
\]

Nếu $a > 1$ thì $a^\alpha > a^\beta$ khi và chỉ khi $\alpha > \beta$.

Nếu $a < 1$ thì $a^\alpha > a^\beta$ khi và chỉ khi $\alpha < \beta$.

\vspace{1em}

\textbf{Ví dụ 1.12} \textit{Rút gọn biểu thức}
\[
E = \frac{a^{\sqrt{7}+1} . a^{2-\sqrt{7}}}{(a^{\sqrt{2}-2})^{\sqrt{2}+2}} \qquad (a > 0).
\]

\textbf{\color{blue}\textit{Giải.}} Với $a > 0$, ta có
\[
E = \frac{a^{\sqrt{7}+1+2-\sqrt{7}}}{a^{(\sqrt{2}-2)(\sqrt{2}+2)}} = \frac{a^3}{a^{-2}} = a^5.
\]

\textbf{Ví dụ 1.13} \textit{Không sử dụng máy tính, hãy so sánh các số $5^{2\sqrt{3}}$ và $5^{3\sqrt{2}}$.}

\textbf{\color{blue}\textit{Giải.}} Ta có $2\sqrt{3} = \sqrt{12}, 3\sqrt{2} = \sqrt{18}$.

Do $12 < 18$ nên $2\sqrt{3} < 3\sqrt{2}$.

Vì cơ số $5$ lớn hơn $1$ nên $5^{2\sqrt{3}} < 5^{3\sqrt{2}}$.
\newpage
\section{HÀM SỐ LŨY THỪA}

\subsection{KHÁI NIỆM}

\textbf{Định nghĩa 2.1} \textit{Ta đã biết các hàm số $y = x^n \ (n \in \mathbb{N}^*), y = \frac{1}{x} = x^{-1}, y = \sqrt{x} = x^{\frac{1}{2}} \ (x > 0)$. Bây giờ, ta xét hàm số $y = x^\alpha$ với $\alpha$ là số thực cho trước.}

\begin{center}
\textit{Hàm số $y = x^\alpha$ với $\alpha \in \mathbb{R}$, được gọi là hàm số lũy thừa.}
\end{center}

\textit{Chẳng hạn, các hàm số $y = x, y = x^2, y = \frac{1}{x^4}, y = x^{\frac{1}{3}}, y = x^{\sqrt{2}}, y = x^\pi$ là những hàm số lũy thừa.}

\vspace{0.5cm}

\textbf{Chú ý 2.2} \textit{Tập xác định của hàm số lũy thừa $y = x^\alpha$ tùy thuộc vào giá trị của $\alpha$. Cụ thể,}

\begin{quote}
\textit{Với $\alpha$ nguyên dương, tập xác định là $\mathbb{R}$;}

\textit{Với $\alpha$ nguyên âm hoặc bằng $0$, tập xác định là $\mathbb{R} \setminus \{0\}$;}

\textit{Với $\alpha$ không nguyên, tập xác định là $(0; +\infty)$.}
\end{quote}
\subsection{ĐẠO HÀM CỦA HÀM SỐ LŨY THỪA}

Một cách tổng quát, người ta chứng minh được hàm số lũy thừa
\[
y=x^\alpha \qquad (\alpha\in\mathbb{R})
\]
có đạo hàm với mọi \(x>0\) và
\[
\left(x^\alpha\right)'=\alpha x^{\alpha-1}.
\]

\textbf{Ví dụ 2.3}
\[
\begin{aligned}
\text{a)}\quad
\left(x^{\frac34}\right)'
&=\frac34x^{\frac34-1}
=\frac{3}{4\sqrt[4]{x}}
\qquad (x>0);\\[6pt]
\text{b)}\quad
\left(x^{\sqrt3}\right)'
&=\sqrt3x^{\sqrt3-1}
\qquad (x>0).
\end{aligned}
\]

\textbf{Chú ý 2.4.}
Công thức tính đạo hàm của hàm hợp đối với hàm số lũy thừa có dạng
\[
\left(u^\alpha\right)'=\alpha u^{\alpha-1}\cdot u'.
\]

\subsection{KHẢO SÁT HÀM SỐ LŨY THỪA \(y=x^\alpha\)}

\begin{samepage}

Tập xác định của hàm số lũy thừa \(y=x^\alpha\) luôn chứa khoảng
\[
(0;+\infty)
\]
với mọi \(\alpha\in\mathbb{R}\). Trong trường hợp tổng quát, ta khảo sát hàm số
\(y=x^\alpha\) trên khoảng này (gọi là tập khảo sát).

\vspace{-5pt}

\begin{center}
\small
\setlength{\tabcolsep}{5pt}
\renewcommand{\arraystretch}{0.9}

\begin{tabular}{@{}p{0.45\textwidth}|p{0.45\textwidth}@{}}

% ===== TIÊU ĐỀ =====
{\centering \(y=x^\alpha,\quad \alpha>0\)\par}
&
{\centering \(y=x^\alpha,\quad \alpha<0\)\par}
\\[8pt]

\hline
\\[-6pt]

% ===== MỤC 1 =====
\textbf{1. Tập khảo sát:} \((0;+\infty)\).
&
\textbf{1. Tập khảo sát:} \((0;+\infty)\).
\\[5pt]


% ===== MỤC 2 =====
\begin{minipage}[t]{\linewidth}
\textbf{2. Sự biến thiên}

\[
y'=\alpha x^{\alpha-1}>0,\quad \forall x>0.
\]

Giới hạn đặc biệt:
\[
\lim_{x\to0^+}x^\alpha=0,\qquad
\lim_{x\to+\infty}x^\alpha=+\infty.
\]

Tiệm cận: Không có.
\end{minipage}
&
\begin{minipage}[t]{\linewidth}
\textbf{2. Sự biến thiên}

\[
y'=\alpha x^{\alpha-1}<0,\quad \forall x>0.
\]

Giới hạn đặc biệt:
\[
\lim_{x\to0^+}x^\alpha=+\infty,\qquad
\lim_{x\to+\infty}x^\alpha=0.
\]

Tiệm cận:

Trục \(Ox\) là tiệm cận ngang,

Trục \(Oy\) là tiệm cận đứng của đồ thị.
\end{minipage}
\\[5pt]

% ===== MỤC 3 =====
\textbf{3. Bảng biến thiên}

\[
\begin{array}{c|cc}
x&0&+\infty\\ \hline
y'&+&\\ \hline
y&0&\nearrow+\infty
\end{array}
\]
&
\textbf{3. Bảng biến thiên}

\[
\begin{array}{c|cc}
x&0&+\infty\\ \hline
y'&-&\\ \hline
y&+\infty&\searrow0
\end{array}
\]
\\[5pt]

% ===== MỤC 4 =====
\textbf{4. Đồ thị (H. 28 với \(\alpha>0\)).}
&
\textbf{4. Đồ thị (H. 28 với \(\alpha<0\)).}

\end{tabular}
\end{center}

\end{samepage}

\begin{center}
    \includegraphics[width=0.6\textwidth]{anh2.png}
\end{center}

Đồ thị của hàm số lũy thừa $y = x^\alpha$ luôn đi qua điểm $(1;1)$.\\
Trên Hình 28 là đô thị của hàm số lũy thừa trên khoảng $(0;+∞)$ ứng với các giá trị khác nhau
của $\alpha$.\\
\textbf{Chú ý 2.5} Khi khảo sát hàm số lũy thừa với số mũ cụ thể, ta phải xét hàm số đó trên toàn bộ tập xác định của nó.
\newpage
\section{MỘT SỐ BÀI TOÁN LIÊN QUAN ĐẾN LŨY THỪA}
\textbf{Bài 1:} Rút gọn biểu thức

\[
P=\frac{\left(a^{\sqrt3-1}\right)^{\sqrt3+1}}
{a^{4-\sqrt5}\cdot a^{\sqrt5-2}}.
\]

\textbf{Giải:}
\[
P=\frac{a^{(\sqrt3-1)(\sqrt3+1)}}
{a^{4-\sqrt5+\sqrt5-2}}
=\frac{a^2}{a^2}=1.
\]\\
\textbf{Bài 2:} Hãy biến đổi và viết lại $P$ dưới dạng một lũy thừa với số mũ hữu tỉ của $x$.
\[
P=\sqrt{x^3}\cdot\sqrt[3]{x^2}\cdot\sqrt[6]{x^5}
\quad (x>0)
\]

\textbf{Giải:}
\[
P=(x^3)^{\frac12}\cdot(x^2)^{\frac13}\cdot(x^5)^{\frac16}
\]
\[
=x^{\frac32}\cdot x^{\frac23}\cdot x^{\frac56}
=x^{\frac32+\frac23+\frac56}
=x^3.
\]\\
\textbf{Bài 3:} Không sử dụng máy tính cầm tay, hãy so sánh các số:

\hspace{1.5cm}\noindent a) $2^{300}$ và $3^{200}$; \hfill b) $\left(\sqrt{5}\right)^{-\frac{2}{3}}$ và $\sqrt[3]{4}$. \hspace*{0.2\textwidth}

\textbf{Giải:}

\noindent a) Ta có $2^{300} = (2^3)^{100} = 8^{100}$; $3^{200} = (3^2)^{100} = 9^{100}$.\\
Do $8 < 9$ và $100 > 0$ nên $8^{100} < 9^{100}$ hay $2^{300} < 3^{200}$.

\vspace{0.3cm}

\noindent b) Ta có: $\sqrt[3]{4} = \sqrt[3]{2^2} = 2^{\frac{2}{3}} = \left(\dfrac{1}{2}\right)^{-\frac{2}{3}}$.

\vspace{0.2cm}

\noindent Do $\sqrt{5} > 1 > \dfrac{1}{2}$ và $-\dfrac{2}{3} < 0$ nên $\left(\sqrt{5}\right)^{-\frac{2}{3}} < \left(\dfrac{1}{2}\right)^{-\frac{2}{3}}$ hay $\left(\sqrt{5}\right)^{-\frac{2}{3}} < \sqrt[3]{4}$.\\

\textbf{Bài 4:} Định luật thứ ba của Kepler về quỹ đạo chuyển động cho biết cách ước tính khoảng thời gian $P$ (tính theo năm Trái Đất) mà một hành tinh cần để hoàn thành một quỹ đạo quay quanh Mặt Trời. Khoảng thời gian đó được xác định bởi hàm số $P = d^{\frac{3}{2}}$, trong đó $d$ là khoảng cách từ hành tinh đó đến Mặt Trời tính theo đơn vị thiên văn $\text{AU}$ ($1\text{ AU}$ là khoảng cách từ Trái Đất đến Mặt Trời, tức là $1\text{ AU}$ khoảng $93\ 000\ 000$ dặm) (\textit{Nguồn: R.I. Charles et al., Algebra 2, Pearson}). Hỏi Sao Hoả quay quanh Mặt Trời thì mất bao nhiêu năm Trái Đất (làm tròn kết quả đến hàng phần trăm)? Biết khoảng cách từ Sao Hoả đến Mặt Trời là $1,52\text{ AU}$.

\vspace{0.3cm}

\textbf{Giải:}

\noindent Vì khoảng cách từ Sao Hoả đến Mặt Trời là $d = 1,52\text{ AU}$, nên khoảng thời gian để Sao Hoả hoàn thành một quỹ đạo quay quanh Mặt Trời là:
\[
P = (1,52)^{\frac{3}{2}} = \sqrt{(1,52)^3} \approx 1,87 \quad \text{(năm Trái Đất)}.
\]

\noindent Vậy Sao Hoả quay quanh Mặt Trời mất khoảng $1,87$ năm Trái Đất.

\textbf{Bài 5:} Một quốc gia ở châu Á có dân số vào năm 2021 là 19 triệu người. Dự kiến sau 30 năm, dân số của quốc gia này sẽ tăng gấp đôi. Công thức ước tính dân số $A$ (đơn vị: triệu người) sau $t$ năm kể từ năm 2021 là:
\[
A = 19 \cdot 2^{\frac{t}{30}}
\]
Hỏi nếu tốc độ tăng trưởng giữ nguyên thì sau 20 năm nữa, dân số quốc gia này đạt khoảng bao nhiêu triệu người? (Làm tròn kết quả đến hàng phần trăm).

\vspace{0.3cm}

\textbf{Giải:}

\noindent Ta có:
\[
A = 19 \cdot 2^{\frac{20}{30}} = 19 \cdot 2^{\frac{2}{3}} = 19 \cdot \sqrt[3]{2^2} \approx 30,16 \quad \text{(triệu người)}.
\]

\noindent Vậy sau 20 năm nữa, dân số của quốc gia đó sẽ vào khoảng $30,16$ triệu người.
\newpage

\section{MỘT SỐ ỨNG DỤNG CỦA LŨY THỪA}
Lũy thừa và àm số lũy thừa được sử dụng rộng rãi trong nhiều lĩnh vực của đời sống, khoa học và kỹ thuật, nhờ vào khả năng mô hình hóa các hiện tượng tăng trưởng hoặc suy giảm theo hàm số mũ.\\
\textbf{1. Ứng dụng trong kinh tế:}\\
Trong kinh tế, hàm số lũy thừa được áp dụng để tính lãi suất kép, một phương pháp tính lãi phổ biến cho các khoản đầu tư dài hạn, giúp dự đoán giá trị tài sản tăng theo thời gian...
\begin{center}
    \includegraphics[width=0.6\textwidth]{anh3.png}
\end{center}
\\
\textbf{2. Ứng dụng trong lĩnh vực đời sống và xã hội:}\\
Hàm số lũy thừa cũng được dùng để mô hình hóa sự tăng trưởng dân số của một khu vực, từ đó có thể dự báo và quy hoạch phát triển....
\begin{center}
    \includegraphics[width=0.6\textwidth]{anh4.png}
\end{center}
\newpage
\textbf{3. Ứng dụng trong lĩnh vực khoa học kỹ thuật:}\\
 Bài toán liên quan đến sự phóng
xạ, tính toán các cơn dư chấn do động đất, cường độ và mức cường độ âm thanh ...
\begin{center}
    \includegraphics[width=0.6\textwidth]{anh5.png}
\end{center}
\textbf{4. Ứng dụng trong sinh học:}\\
 Trong sinh học, hàm số lũy thừa được dùng để mô tả sự tăng trưởng của vi khuẩn trong một môi trường nuôi cấy. Ví dụ, nếu bắt đầu với 100 vi khuẩn, sau một thời gian nhất định, số lượng vi khuẩn có thể tăng lên đáng kể, tuân theo một hàm số lũy thừa.
\begin{center}
    \includegraphics[width=0.6\textwidth]{anh6.png}
\end{center}
\textbf{5. Ứng dụng trong y học:}\\
 Trong y học, các hàm lũy thừa được dùng để mô tả sự phân rã của thuốc trong cơ thể, giúp xác định lượng thuốc còn lại sau một khoảng thời gian nhất định.
 \newpage
 \textbf{TÀI LIỆU THAM KHẢO}\\
1. Trần Văn Hạo(Tổng Chủ biên)-Vũ Tuấn (Chủ biên), Lê Thị Thiên Hương, Nguyễn Tiến Tài,
Cấn Văn Tuất, Toán 12 giải tích, Nhà xuất bản giáo dục Việt Nam, (2008). \\
2. Các dạng bài tập về hàm lũy thừa: https://vietjack.me/cac-dang-bai-tap-ve-cong-thuc-luy-thua-logarit-toan-lop
12-31137.html [3] Ứng dụng của hàm số lũy thừa: https://toanmath.com/2017/11/ung-dung ham-so-luy-thua-ham-so-mu-va-ham-so-logarit-de-giai-cac-bai-toan-thuc-te-lien-quan.html\\
3. Ứng dụng thực tế của hàm số lũy thừa https://rdsic.edu.vn/blog/toan/phan-tich-ham-so-luy-thua-12-cac-dac-tinh-va-ung-dung-ban-can-biet-vi-cb.html





\end{document}